% ==============================================================================
% preamble_template_quickref.tex
% มาตรฐานการตั้งค่า XeLaTeX ตำราวิชาการภาษาไทย ขนาด 6 x 9 นิ้ว
% อ้างอิงมาตรฐาน: thai-academic-latex-formatting
% พัฒนาโดย: ผศ.ดร.ชีวะ ทัศนา (มหาวิทยาลัยราชภัฏรำไพพรรณี)
% ==============================================================================

\documentclass[11pt,twoside,openright]{book}

% ------------------------------------------------------------------------------
% 1. ขนาดรูปเล่มและระยะขอบ (Physical Geometry: 6 x 9 นิ้ว)
% ------------------------------------------------------------------------------
\usepackage{geometry}
\geometry{
    paperwidth=6in,
    paperheight=9in,
    inner=1.0in,     % ขอบในสำหรับเข้าเล่ม (Binding Gutter)
    outer=0.75in,    % ขอบนอก
    top=0.85in,      % ขอบบน
    bottom=0.85in,   % ขอบล่าง
    headheight=14pt,
    headsep=18pt,
    footskip=24pt
}

% ------------------------------------------------------------------------------
% 2. ฟอนต์และการจัดพิมพ์ภาษาไทย (Thai Typography & Micro-spacing)
% ------------------------------------------------------------------------------
\usepackage{fontspec}
\usepackage{setspace}
\usepackage{indentfirst} % บังคับย่อหน้าแรกหลังหัวข้อเสมอ

% กำหนดฟอนต์หลัก (Sarabun จากโฟลเดอร์ fonts/ หรือระบบ)
\setmainfont{Sarabun}[
    Path = fonts/,
    Extension = .ttf,
    UprightFont = *-Regular,
    BoldFont = *-Bold,
    ItalicFont = *-Italic,
    BoldItalicFont = *-BoldItalic
]

% ระบบตัดคำภาษาไทยอัตโนมัติใน XeLaTeX
\XeTeXlinebreaklocale "th"
\XeTeXlinebreakskip = 0pt plus 1pt

% ระยะห่างบรรทัด 1.18 เท่าเพื่อความสบายตาและรองรับวรรณยุกต์ 4 ชั้น
\setstretch{1.18}

% ย่อหน้า 3em ตามความต้องการของผู้ใช้งาน ไม่เว้นวรรคระหว่างย่อหน้า
\setlength{\parindent}{3em}
\setlength{\parskip}{0pt}

% ------------------------------------------------------------------------------
% 3. แพ็กเกจสนับสนุนงานวิชาการและกราฟิก
% ------------------------------------------------------------------------------
\usepackage{graphicx}
\usepackage[dvipsnames,table]{xcolor}
\usepackage{tikz}
\usepackage{tcolorbox}
\tcbuselibrary{skins,breakable}
\usepackage{booktabs}
\usepackage{tabularx}
\usepackage{amsmath,amssymb}
\usepackage{needspace}
\usepackage{fancyhdr}
\usepackage{caption}

% ------------------------------------------------------------------------------
% 4. รูปแบบคำบรรยายภาพและตาราง (Captions: ภาพที่ / ตารางที่ ไร้โคลอน)
% ------------------------------------------------------------------------------
\DeclareCaptionLabelSeparator{space}{ }
\captionsetup{
    labelsep=space,
    labelfont=bf,
    font=small
}
% ตาราง: จัดชิดซ้ายทั้งเล่ม
\captionsetup[table]{
    singlelinecheck=false,
    justification=raggedright,
    position=top
}
% ภาพ: กึ่งกลาง
\captionsetup[figure]{
    justification=centering,
    position=bottom
}
\renewcommand{\figurename}{ภาพที่}
\renewcommand{\tablename}{ตารางที่}

% ------------------------------------------------------------------------------
% 5. หัวกระดาษและท้ายกระดาษ (Running Headers: หน้าคู่ / หน้าคี่)
% ------------------------------------------------------------------------------
\pagestyle{fancy}
\fancyhf{}
\fancyhead[LE]{\small\thepage\quad\textbar\quad\textsf{คู่มือปฏิบัติการ Digital Agri xAI}}
\fancyhead[RO]{\small\textsf{\leftmark}\quad\textbar\quad\thepage}
\renewcommand{\headrulewidth}{0.4pt}
\renewcommand{\chaptermark}[1]{\markboth{บทที่ \thechapter\ #1}{}}

% ------------------------------------------------------------------------------
% 6. หัวข้อและการป้องกันหัวข้อตกขอบ (Headings with \needspace)
% ------------------------------------------------------------------------------
\usepackage{titlesec}

% ปรับแต่งรูปแบบบท (Chapter Heading: ชิดขวา ตัวหนาเด่นชัด)
\titleformat{\chapter}[display]
    {\normalfont\huge\bfseries\raggedleft}
    {\fontsize{18pt}{22pt}\selectfont\color{ForestGreen}บทที่ \thechapter}
    {12pt}
    {\fontsize{22pt}{28pt}\selectfont\color{black!90}}

\titlespacing*{\chapter}{0pt}{20pt}{35pt}

% ก่อนขึ้น Section ป้องกันหัวข้อโดดเดี่ยวท้ายหน้า
\newcommand{\safeSection}[1]{%
    \needspace{5\baselineskip}%
    \section{#1}%
}
\newcommand{\safeSubsection}[1]{%
    \needspace{4\baselineskip}%
    \subsection{#1}%
}

% ------------------------------------------------------------------------------
% 7. เอกสารอ้างอิงมาตรฐาน APA 7th Edition
% ------------------------------------------------------------------------------
\newcommand{\bibentry}[1]{\noindent\hangindent=2.5em\hangafter=1 #1\par\vspace{7pt}}

% ------------------------------------------------------------------------------
% 8. กล่องข้อความวิชาการ (Academic Callout Box)
% ------------------------------------------------------------------------------
\newtcolorbox{academicNote}[1][]{
    enhanced,
    breakable,
    colback=Emerald!6!white,
    colframe=Emerald!75!black,
    fonttitle=\bfseries,
    title=💡 ข้อสังเกตและแนวทางปฏิบัติ,
    arc=4mm,
    boxrule=0.8pt,
    left=15pt, right=15pt, top=12pt, bottom=12pt,
    #1
}
